%% P10 --- SVD, low-rank approximation and conditioning
%%
%% Stub. The brief below is the contract for this program: what it must argue,
%% and what it must buy the reader. Writing the program means deleting the
%% \programstub{} block and replacing it with the program. If the brief turns
%% out to be wrong once the mathematics has been worked, change it and record
%% the contradiction in CLAUDE.md under *Resolved questions*.

\program{SVD, low-rank approximation and conditioning}
\label{prog:P10}

\programstub{%
Argument: every matrix, square or not, factors into rotate--stretch--rotate.
Truncating the stretch gives the provably best approximation of a given rank
(Eckart--Young, stated). The ratio of largest to smallest singular value is
the condition number, and it is the amplification factor from input error to
output error. Payoff: the singular-value spectrum of a real embedding
matrix, plotted, showing how few directions carry the energy --- the
empirical case for LoRA and for embedding compression, measured; the
pseudoinverse as least squares done properly; why the normal equations
square the condition number and a QR solve does not, which is the concrete
form of \enquote{do not invert}; and the condition number as the number that
predicts how many iterations an optimiser will need, handed forward to P19.

\medskip
\textbf{Planned length:} 60 frames.
}
